\documentclass[11pt,a4paper]{article}
\usepackage[T1]{fontenc}
\usepackage[utf8]{inputenc}
\usepackage{lmodern}
\usepackage{amsmath,amssymb,booktabs}
\usepackage[hidelinks]{hyperref}
\title{Privacy-Preserving Facial Behaviour Estimation on Commodity RGB Cameras}
\author{}
\date{}
\begin{document}
\maketitle
\begin{abstract}
What performance and system costs remain when facial pixels are never persisted? The manuscript defines a reproducible research design that tests an explicit methodological contribution against appropriate non-adaptive or low-capacity baselines. All comparisons are prospective; no empirical improvement or validated model performance is claimed at this stage.
\end{abstract}
\section{Introduction and related methodological problem}
What performance and system costs remain when facial pixels are never persisted? The scientific goal is to distinguish a genuine computational contribution from gains due to subject overlap, dataset bias, confounded acquisition settings, or favourable synthetic conditions.
\section{Research objectives and hypotheses}
Compare ephemeral landmarks-only pipelines with on-device image inference for privacy/utility/latency trade-offs. Hypothesis H1 predicts better pre-specified robustness or accuracy than an appropriate baseline under independently held-out evaluation; hypothesis H2 requires that the measured gain remains interpretable after accounting for uncertainty and deployment constraints. Both hypotheses are proposals, not tested findings.
\section{Materials and data eligibility}
Licensed Aff-Wild2 or AFEW-VA annotations after access; optional local non-recorded webcam N=1 for deployment timing. Data usage requires source licensing, stable event/label definitions, and a documented eligibility gate. Any unavailable or incomparable source must be excluded with an explicit reason.
\section{Methods and experimental design}
Separate facial AU/landmark estimation from affect target inference; frame-private flow; compare landmarks-only, compact video model and no-video baseline; benchmark asynchrony, illumination and occlusion. Explicitly avoid face identity recognition. Hyperparameters and calibration are selected only from training splits. Evaluation units and confidence intervals must correspond to independent subjects, sessions, devices or scenarios rather than artificially replicated samples.
\section{Statistical analysis and primary outcomes}
CCC valence/arousal where labels valid, AU F1, p95 latency, RAM, energy when measurable, retention audit, privacy design verification. Every primary comparison will report uncertainty intervals, failure cases, negative controls and the actual sample/split sizes. Planned sensitivity analyses vary missingness, corruption and acquisition conditions.
\section{Results}
No empirical dataset analysis or confirmatory comparison has been completed for this manuscript. This section will report results only after authorised data collection or reproducible benchmark execution; planned metrics are not observed effects.
\section{Discussion, limitations and continuation criteria}
Facial expressivity is not emotional ground truth; landmark streams can themselves be identifying; dataset use restrictions. Verify all frame persistence disabled in default path, scientifically meaningful baseline and explicit facial proxy limitations. A novel method cannot be asserted solely from a different application domain; comparison against current prior art and reproducibility on an independent setting are required.
\section{Reproducibility and provenance}
The original research plan is preserved in \texttt{source\_plan.md}; repository and exact source commit are recorded in \texttt{provenance.json}. No sensitive recordings or private credentials are embedded.
\end{document}
